\section{Introduction}

Suppose two firms each have an even chance of a return gain. Their returns
move together in one regime and in opposite directions in another. Exact
marginal predictions are compatible with either dependence structure.
What would a tabular prediction task have to preserve for its learned
features to measure financial proximity?
\lean{FreeFoundations.Alignment.marginals_free_of_state}
\lean{FreeFoundations.Alignment.joint_depends_on_state}

For firms with loadings $b_i,b_j$ on centered common shocks $q$, define
proximity through the root mean square of $(b_i-b_j)^\top q$. This gives
distance an economic unit, the size of the difference in the firms'
return responses to the same shocks. A training task guarantees recovery
of a history-based exposure precisely when its ideal prediction separates
histories with different exposures. Even then, the choice of metric
determines whose returns enter each firm's peer average.
\lean{FreeFoundations.Score.response_sq}

We compare a frozen encoder with the characteristic-history distance of
\citet{fernandez2011} and with peers selected directly from the same filings.

In our US reimplementation of Fern\'andez's spatial CAPM, the joint
characteristic kernel spreads weight across so many issuers that firms with
quite different histories receive much the same peer return. To separate the
choice of distance from the concentration of the weights, we compare
distances under a common ten-neighbour rule and then examine how the
comparison changes as weight extends to more distant firms.

Across the matched quarterly panel, peers selected by TabPFN from four
financial ratios improve return fit over Fern\'andez's joint distance even
after adjustment for the five Fama--French factors and momentum. To
determine how much of that fit is already available in the financial inputs,
we compare each encoder with a graph built directly from the features it
receives. Under market controls, the estimated TabPFN gain on the four
ratios is close to zero, with an interval spanning gains and losses. When
both methods draw on a broad panel of financial disclosures, TabPFN yields
the largest estimated gain over its direct benchmark among the twelve
combinations we examine.

The economic interpretation of these comparisons turns on a distinction
between recovering a firm's exposure and identifying its peers. Two
representations can preserve exactly the same exposure while placing
different firms nearest to it. We show how such a change of coordinates
can alter the network and bound the change in implied covariance at a
common spatial coefficient. In the selected broad-input model we also
measure how many peers remain when we discard all but the leading
valuation coordinate. This measures the cost of using a one-dimensional
summary for peer selection.
